%% Phan cua Thanh vien 2
\subsection{Phát biểu bài toán}\label{sec:problem}

\begin{definition}[Bandwidth Coloring]\label{def:bcp}
Cho đồ thị vô hướng $G=(V,E)$ với $n=|V|$ và trọng số nguyên dương $d_{uv}$ trên mỗi cạnh. Một \emph{phép tô hợp lệ} là hàm $c:V\to\mathbb{Z}_{>0}$ thỏa mãn
\begin{equation}
|c(u)-c(v)|\ge d_{uv}\qquad\forall (u,v)\in E. \label{eq:sep}
\end{equation}
Bài toán yêu cầu tìm phép tô hợp lệ có màu lớn nhất (gọi là \emph{span}) nhỏ nhất:
\begin{equation}
\chi_B(G)=\min_{c}\ \max_{v\in V}c(v). \label{eq:obj}
\end{equation}
\end{definition}

Hình~\ref{fig:example} là một ví dụ nhỏ. Đỉnh $v_1$ và $v_4$ không kề nhau nên có thể dùng chung màu. Không có phép tô hợp lệ nào có span 5, nên $\chi_B(G)=6$.

\begin{figure}[pos=t]
\centering
\begin{tikzpicture}[vert/.style={circle,draw,minimum size=6.2mm,inner sep=0pt,font=\footnotesize},
  peak/.style={vert,draw=blue!70!black,fill=blue!10,thick}, wt/.style={font=\footnotesize}]
\node[vert] (v0) at (0,0) {$v_1$};
\node[vert] (v1) at (1.5,1.25) {$v_2$};
\node[vert] (v2) at (1.5,-1.25) {$v_3$};
\node[vert] (v3) at (3,0) {$v_4$};
\draw (v0)--node[wt,above left]{3}(v1);
\draw (v0)--node[wt,below left]{2}(v2);
\draw (v1)--node[wt,right]{2}(v2);
\draw (v1)--node[wt,above right]{1}(v3);
\draw (v2)--node[wt,below right]{3}(v3);
\node[font=\footnotesize] at (1.5,-2.15) {(a) đồ thị $G$, số trên cạnh là $d_{uv}$};
\begin{scope}[shift={(4.6,0.55)}]
\draw (0,0)--(7.6,0);
\draw (0.3,-0.06)--(0.3,0.06); \node[font=\scriptsize,below] at (0.3,-0.04) {1};
\draw (1.3,-0.06)--(1.3,0.06); \node[font=\scriptsize,below] at (1.3,-0.04) {2};
\draw (2.3,-0.06)--(2.3,0.06); \node[font=\scriptsize,below] at (2.3,-0.04) {3};
\draw (3.3,-0.06)--(3.3,0.06); \node[font=\scriptsize,below] at (3.3,-0.04) {4};
\draw (4.3,-0.06)--(4.3,0.06); \node[font=\scriptsize,below] at (4.3,-0.04) {5};
\draw (5.3,-0.06)--(5.3,0.06); \node[font=\scriptsize,below] at (5.3,-0.04) {6};
\draw (6.3,-0.06)--(6.3,0.06); \node[font=\scriptsize,below] at (6.3,-0.04) {7};
\draw (7.3,-0.06)--(7.3,0.06); \node[font=\scriptsize,below] at (7.3,-0.04) {8};
\node[vert] at (0.3,0.42) {$v_1$};
\node[vert] at (3.3,0.42) {$v_2$};
\node[peak] at (7.3,0.42) {$v_3$};
\node[vert] at (0.3,1.10) {$v_4$};
\node[font=\footnotesize,anchor=west] at (1,1.1) {(b) phép tô khả thi, span 8};
\end{scope}
\begin{scope}[shift={(4.6,-1.95)}]
\draw (0,0)--(7.6,0);
\draw (0.3,-0.06)--(0.3,0.06); \node[font=\scriptsize,below] at (0.3,-0.04) {1};
\draw (1.3,-0.06)--(1.3,0.06); \node[font=\scriptsize,below] at (1.3,-0.04) {2};
\draw (2.3,-0.06)--(2.3,0.06); \node[font=\scriptsize,below] at (2.3,-0.04) {3};
\draw (3.3,-0.06)--(3.3,0.06); \node[font=\scriptsize,below] at (3.3,-0.04) {4};
\draw (4.3,-0.06)--(4.3,0.06); \node[font=\scriptsize,below] at (4.3,-0.04) {5};
\draw (5.3,-0.06)--(5.3,0.06); \node[font=\scriptsize,below] at (5.3,-0.04) {6};
\draw (6.3,-0.06)--(6.3,0.06); \node[font=\scriptsize,below] at (6.3,-0.04) {7};
\draw (7.3,-0.06)--(7.3,0.06); \node[font=\scriptsize,below] at (7.3,-0.04) {8};
\node[vert] at (3.3,0.42) {$v_1$};
\node[vert] at (0.3,0.42) {$v_2$};
\node[peak] at (5.3,0.42) {$v_3$};
\node[vert] at (1.3,0.42) {$v_4$};
\node[font=\footnotesize,anchor=west] at (1,1.1) {(c) phép tô đọc từ mô hình SAT, span 6};
\end{scope}
\end{tikzpicture}

\caption{Ví dụ BCP với 4 đỉnh và 5 cạnh. Đỉnh tô đậm là đỉnh có màu lớn nhất.}
\label{fig:example}
\end{figure}
